%! The Rational Parent Function \& Transformations — Unit 4, Lesson 4.4 — Group Activity (Answer Key)
\documentclass[10pt]{article}
\usepackage{algebra2-article}
\usepackage{algebra2-key}   % pulls in -boxes; do NOT also load -boxes
\newcommand{\TallMath}[1]{$\displaystyle #1\rule[-1.4em]{0pt}{3.2em}$}
% Standard coordinate grid + axes: {xmin}{xmax}{ymin}{ymax}
\newcommand{\lgrid}[4]{%
  \draw[step=1, linegray!40, very thin] (#1,#3) grid (#2,#4);
  \draw[->, semithick, charcoal] (#1,0)--(#2,0) node[below right, font=\scriptsize]{$x$};
  \draw[->, semithick, charcoal] (0,#3)--(0,#4) node[left, font=\scriptsize]{$y$};}

\begin{document}

\pageheader{Unit 4, Lesson 4.4}{Group Activity --- Answer Key}
\namepartnerperiod

\vspace{0.08in}

\begin{headlinebox}{forestbg}
\small\textbf{Move the Walls.} One habit carries everything today: \emph{find the two asymptotes
first.} They are the parent's axes, relocated --- $x=h$ from the denominator and $y=k$ from the number
tacked on outside. Once those two dashed lines are down, the branches have nowhere else to go. And keep
Unit 1's warning in front of you: the number \emph{inside} does the opposite of what it says.
\end{headlinebox}

\vspace{0.06in}

\begin{tcolorbox}[colback=white, colframe=black!40, title=\textbf{Tier R --- Remediate},
    fonttitle=\bfseries, arc=1.5mm, left=3mm, right=3mm, top=2mm, bottom=2mm, breakable]
\textbf{Part 1 --- One move at a time.} For each function, name both asymptotes and say what the parent
$y=\frac1x$ did.
\begin{enumerate}[leftmargin=1.9em, itemsep=8pt]
  \item $y=\dfrac{1}{x-5}$ \\[3pt]
        VA: \ans{$x=5$} \quad HA: \ans{$y=0$} \quad The parent moved \ans{right $5$}
  \item $y=\dfrac{1}{x}-4$ \\[3pt]
        VA: \ans{$x=0$} \quad HA: \ans{$y=-4$} \quad The parent moved \ans{down $4$}
  \item $y=\dfrac{1}{x+6}$ \\[3pt]
        VA: \ans{$x=-6$} \quad HA: \ans{$y=0$} \quad The parent moved \ans{left $6$}
  \item $y=\dfrac{2}{x-1}+3$ \\[3pt]
        $a=$ \ans{$2$} \quad $h=$ \ans{$1$} \quad $k=$ \ans{$3$} \quad
        Center: $($\ans{$1$}$,\,$\ans{$3$}$)$ \\[3pt]
        VA: \ans{$x=1$} \quad HA: \ans{$y=3$} \quad Domain: \ans{$x\neq1$} \quad
        Range: \ans{$y\neq3$}
\end{enumerate}

\vspace{5pt}
\textbf{Part 2 --- Match the graph to the equation.} Write \textbf{A}, \textbf{B}, or \textbf{C} in each
blank. The dashed navy lines are the asymptotes.
\begin{center}
\begin{tikzpicture}[scale=0.30]
  % All three windows share y from -5 to 5 so the graphs sit on a common baseline.
  \begin{scope}
    \lgrid{-3}{6}{-5}{5}
    \draw[navy, dashed, semithick] (2,-5)--(2,5);
    \begin{scope} \clip (-3,-5) rectangle (6,5);
      \draw[very thick, forest, domain=-3:1.8, samples=100, smooth] plot (\x,{1/((\x)-2)});
      \draw[very thick, forest, domain=2.2:6, samples=100, smooth] plot (\x,{1/((\x)-2)});
    \end{scope}
    \node[charcoal, font=\small\bfseries] at (1.5,-6.8){A};
  \end{scope}
  \begin{scope}[xshift=12cm]
    \lgrid{-5}{5}{-5}{5}
    \draw[navy, dashed, semithick] (-5,3)--(5,3);
    \begin{scope} \clip (-5,-5) rectangle (5,5);
      \draw[very thick, forest, domain=-5:-0.125, samples=110, smooth] plot (\x,{1/(\x)+3});
      \draw[very thick, forest, domain=0.5:5, samples=110, smooth] plot (\x,{1/(\x)+3});
    \end{scope}
    \node[charcoal, font=\small\bfseries] at (0,-6.8){B};
  \end{scope}
  \begin{scope}[xshift=24cm]
    \lgrid{-6}{4}{-5}{5}
    \draw[navy, dashed, semithick] (-2,-5)--(-2,5);
    \draw[navy, dashed, semithick] (-6,-3)--(4,-3);
    \begin{scope} \clip (-6,-5) rectangle (4,5);
      \draw[very thick, forest, domain=-6:-2.5, samples=110, smooth] plot (\x,{1/((\x)+2)-3});
      \draw[very thick, forest, domain=-1.125:4, samples=110, smooth] plot (\x,{1/((\x)+2)-3});
    \end{scope}
    \node[charcoal, font=\small\bfseries] at (-1,-6.8){C};
  \end{scope}
\end{tikzpicture}
\end{center}
\vspace{2pt}
\ans{B} \; $y=\dfrac{1}{x}+3$ \hfill
\ans{C} \; $y=\dfrac{1}{x+2}-3$ \hfill
\ans{A} \; $y=\dfrac{1}{x-2}$ \hfill\hfill

\vspace{5pt}
\textbf{Part 3 --- The parent, from memory.} No graph, no calculator.
\begin{center}
$y=\dfrac1x$: \quad VA \ans{$x=0$} \quad HA \ans{$y=0$} \quad $x$-int \ans{none} \quad
$y$-int \ans{none} \quad symmetry \ans{origin}
\end{center}
Why does this parent have no intercepts, when every other parent you have graphed passes through the
origin? \ansline{A fraction equals $0$ only when its numerator does, and the numerator is $1$; and $x=0$ is excluded from the domain, so there is no output to plot on the $y$-axis.}
\end{tcolorbox}

\begin{tcolorbox}[colback=white, colframe=black!40, title=\textbf{Tier A --- Approaching Proficiency},
    fonttitle=\bfseries, arc=1.5mm, left=3mm, right=3mm, top=2mm, bottom=2mm, breakable]
\textbf{Part 1 --- The full table.} Fill in every cell. Watch the signs.
\begin{center}
\renewcommand{\arraystretch}{1.7}
\begin{tabularx}{\linewidth}{@{}l X X X X@{}}
\hline
\textbf{Function} & \textbf{$a$, $h$, $k$} & \textbf{VA \& HA} & \textbf{Domain \& range} & \textbf{$y$-intercept} \\
\hline
\TallMath{y=\frac{4}{x-1}}      & \ans{$4,\,1,\,0$}   & \ans{$x=1$, $y=0$}   & \ans{$x\neq1$, $y\neq0$}  & \ans{$(0,-4)$} \\
\TallMath{y=\frac{1}{x}-5}      & \ans{$1,\,0,\,-5$}  & \ans{$x=0$, $y=-5$}  & \ans{$x\neq0$, $y\neq-5$} & \ans{none} \\
\TallMath{y=\frac{-2}{x+3}+1}   & \ans{$-2,\,-3,\,1$} & \ans{$x=-3$, $y=1$}  & \ans{$x\neq-3$, $y\neq1$} & \ans{$(0,\frac13)$} \\
\TallMath{y=\frac{6}{x+2}-2}    & \ans{$6,\,-2,\,-2$} & \ans{$x=-2$, $y=-2$} & \ans{$x\neq-2$, $y\neq-2$}& \ans{$(0,1)$} \\
\hline
\end{tabularx}
\end{center}
One of those four functions has \emph{no} $y$-intercept. Which one, and why not? \ans{$y=\frac1x-5$ --- its vertical asymptote \emph{is} the $y$-axis, so $x=0$ is not in the domain}

\vspace{5pt}
\textbf{Part 2 --- Write the equation from the graph.} The dashed navy lines are the asymptotes, and two
points are marked.
\begin{center}
\begin{minipage}{0.48\linewidth}\centering
\begin{tikzpicture}[scale=0.38]
  \lgrid{-4}{7}{-1}{7}
  \draw[navy, dashed, semithick] (1,-1)--(1,7);
  \draw[navy, dashed, semithick] (-4,3)--(7,3);
  \node[navy, font=\scriptsize, above right] at (1,6.2) {$x=1$};
  \node[navy, font=\scriptsize, below right] at (5.4,3) {$y=3$};
  \begin{scope}
    \clip (-4,-1) rectangle (7,7);
    \draw[very thick, forest, domain=-4:0.6, samples=110, smooth] plot (\x,{-2/((\x)-1)+3});
    \draw[very thick, forest, domain=1.4:7, samples=110, smooth] plot (\x,{-2/((\x)-1)+3});
  \end{scope}
  \fill[charcoal] (0,5) circle (3pt) node[above left, font=\scriptsize]{$(0,5)$};
  \fill[charcoal] (3,2) circle (3pt) node[below right, font=\scriptsize]{$(3,2)$};
\end{tikzpicture}
\end{minipage}%
\begin{minipage}{0.48\linewidth}
\small
(a) $h=$ \ans{$1$} \quad $k=$ \ans{$3$}

\vspace{4pt}
(b) Substitute the point $(3,2)$ into $y=\dfrac{a}{x-h}+k$ and solve for $a$:

\ans{$2=\dfrac{a}{3-1}+3 \;\Rightarrow\; \dfrac{a}{2}=-1 \;\Rightarrow\; a=-2$}

$a=$ \ans{$-2$}

\vspace{4pt}
(c) Equation: \ans{$y=\frac{-2}{x-1}+3$}

\vspace{4pt}
(d) Verify with the \emph{other} point, $(0,5)$: \ans{$\frac{-2}{-1}+3=5$ \checkmark}
\end{minipage}
\end{center}

\vspace{4pt}
\textbf{Part 3 --- Error analysis.} A student writes:
\begin{center}
\fbox{\parbox{0.74\linewidth}{\centering \vspace{2pt}
For $y=\dfrac{1}{x+4}-1$: \; ``VA is $x=4$ and HA is $y=1$.'' \\[3pt]
{\small ``I just read the two numbers off the equation.''}\vspace{2pt}}}
\end{center}
\begin{enumerate}[label=(\alph*), leftmargin=1.9em, itemsep=6pt]
  \item Both are wrong. Give the correct asymptotes: VA \ans{$x=-4$} \quad HA \ans{$y=-1$}
  \item Disprove the vertical asymptote with a number: evaluate the function at $x=4$.
        $y=$ \ans{$-\frac78$} \; --- a value exists there, so $x=4$ cannot be a wall.
  \item Now evaluate at $x=-4$: \ans{undefined ($\frac10$)}
  \item In one sentence, state the rule the student broke.
        \ansline{The number inside the denominator does the opposite of what it says --- $x+4$ shifts left, so the wall is at $x=-4$ --- and the number outside is the horizontal asymptote exactly as written, $y=-1$.}
\end{enumerate}
\end{tcolorbox}

\begin{tcolorbox}[colback=white, colframe=black!40, title=\textbf{Tier E --- Extension},
    fonttitle=\bfseries, arc=1.5mm, left=3mm, right=3mm, top=2mm, bottom=2mm, breakable]
\textbf{Part 1 --- Unmask the disguise.} Each function below \emph{is} a transformed $\frac1x$, but it is
not written that way. Rewrite the numerator using the denominator, then read off the asymptotes.
\begin{enumerate}[label=(\alph*), leftmargin=1.9em, itemsep=8pt]
  \item $y=\dfrac{3x+5}{x+1}=\dfrac{3({\color{keyred}\mathbf{x+1}})+{\color{keyred}\mathbf{2}}}{x+1}=$ \ans{$3+\frac{2}{x+1}$}
        \\[3pt]
        $a=$ \ans{$2$} \quad $h=$ \ans{$-1$} \quad $k=$ \ans{$3$} \quad
        VA \ans{$x=-1$} \quad HA \ans{$y=3$}
  \item $y=\dfrac{2x-7}{x-4}=$ \ans{$2+\frac{1}{x-4}$} \\[3pt]
        VA \ans{$x=4$} \quad HA \ans{$y=2$}
  \item Check (b) against Lesson 4.0's degree comparison: the degrees of top and bottom are
        \ans{equal (both $1$)}, so the HA is the ratio of the \ans{leading coefficients}, which is \ans{$\frac21=2$}. Same
        answer? \ans{yes}
\end{enumerate}

\vspace{5pt}
\textbf{Part 2 --- Build one backwards.} Write the equation of a rational function whose vertical
asymptote is $x=-3$, whose horizontal asymptote is $y=4$, and whose graph passes through $(-2,1)$.
\begin{center}
$y=$ \ans{$\frac{-3}{x+3}+4$}
\end{center}
\begin{enumerate}[label=(\alph*), leftmargin=1.9em, itemsep=5pt]
  \item Now write your answer as a \emph{single} fraction (Lesson 4.3, run forwards): \ans{$\frac{4x+9}{x+3}$}
  \item Confirm the horizontal asymptote from that single fraction using the degree comparison.
        \ans{equal degrees $\Rightarrow y=\frac41=4$ \checkmark}
  \item Could a function have those two asymptotes and pass through $(-3,1)$? Explain.
        \ansline{No. $x=-3$ is the vertical asymptote, so $-3$ is excluded from the domain --- there is no point on the graph with that $x$-value at all.}
\end{enumerate}

\vspace{5pt}
\textbf{Part 3 --- Diluting the brine.} A tank holds $20$ grams of salt dissolved in $5$ liters of
water. You pour in $x$ more liters of \emph{pure} water. The concentration is
\[ C(x)=\frac{20}{5+x} \quad \text{grams per liter.} \]
\begin{enumerate}[label=(\alph*), leftmargin=1.9em, itemsep=6pt]
  \item Fill in: $C(0)=$ \ans{$4$}, \; $C(5)=$ \ans{$2$}, \; $C(15)=$ \ans{$1$}, \;
        $C(35)=$ \ans{$0.5$} g/L.
  \item Written as $\dfrac{a}{x-h}+k$: \; $a=$ \ans{$20$}, \; $h=$ \ans{$-5$}, \;
        $k=$ \ans{$0$}. So the horizontal asymptote is \ans{$y=0$}.
  \item What does that horizontal asymptote mean about the salt? \ansline{The concentration drops toward $0$ but never reaches it: the $20$ g of salt never leaves the tank, so no amount of added water makes the brine truly pure.}
  \item The algebra puts a vertical asymptote at $x=$ \ans{$-5$}. Does that number mean anything
        for this tank? \ans{no} \; What values of $x$ actually make sense? \ans{$x\ge0$}
  \item How much water gets the concentration to $2$ g/L? \ans{$5$} L. \; To $1$ g/L?
        \ans{$15$} L. \; To $0.5$ g/L? \ans{$35$} L. \\[3pt]
        So each \emph{halving} of the concentration costs \ans{twice as much} water than the halving before
        it. Explain what the graph is doing near its horizontal asymptote. \ansline{The curve flattens as it approaches $y=0$, so it takes a bigger and bigger run of added water to buy the same drop in concentration --- the graph gets closer to the asymptote but never arrives.}
\end{enumerate}
\end{tcolorbox}

\begin{teachernote}
\textbf{Tier R} is one transformation at a time until item 4. Part 2's matching is the item to watch: a
student who picks \textbf{A} for $y=\frac1x+3$ has the two roles crossed. Part 3 is not a throwaway ---
``no intercepts'' is the fact students most refuse to believe.

\vspace{3pt}
\textbf{Tier A Part 1} row 2 is the point of the table: $y=\frac1x-5$ has \emph{no} $y$-intercept,
because its vertical asymptote is the $y$-axis. Part 2 is A2.F.1b, and step (d) --- verify on a second
point --- is the required habit. In Part 3, beat the student's rule with arithmetic, not a correction
($x=4$ gives an ordinary value, so it cannot be a wall).

\vspace{3pt}
\textbf{Tier E Part 1} is 4.5's algebra a day early and the proof of 4.0's degree rule here: the
leftover constant \emph{is} the ratio of the leading coefficients. Early finishers: $y=\frac{x}{x-2}$.
\textbf{Part 3(d)} is the modeling beat (A2.F.2f) --- $x=-5$ is real mathematics and meaningless
chemistry.
\end{teachernote}

\end{document}
